\section{Conclusion}
\label{sec:conclusion}
Android notification coverage depends on source, lifecycle and process state. Deferred callbacks, frozen intervals and lost observer coverage require different handling. Snapshot certification additionally needs consistent anchoring and reset-safe evidence; it does not certify a later effect. Unsupported stock paths remain \Unknown. Direct refresh and input memoization win the measured comparisons in the client\ifPhysHost; through a host helper on the phone, batched \stamp beat batched refresh only by re-reading fewer facts, with the optimistic policy that certified stale values in the race\fi. On three rented phone models (Pixel~8 and Galaxy~S25 on Android~15, Pixel~9 on Android~16) the measured emulator findings recurred, while foreground callbacks were slower on the Pixel phones (\PhysDeployFgMedian and \PhysNineFgMedian against \FollowAFg) and, on the Pixel~8, the check-to-effect interval longer and file mtimes finer\ifPhysSamsung. On the S25, Samsung's Contacts provider flags its notifications so that a background client hears them without the 10\,s deferral, and Android's freezer never froze a cached client, though Samsung sometimes moved one into its own freezer group after screen off, so the gap classes a client meets depend on the device\fi. The result is an evidence map for choosing a strategy by its deployment, contract and cost.
